% !TeX root = ../main.tex
\section{Introduction}
\label{sec:introduction}

The energy cost of moving data between memory and arithmetic units motivates compute-in-memory (CIM) and processing-in-memory (PIM) architectures for edge inference \citep{horowitz2014computing,peng2020dnnneurosim}. In these architectures, dense matrix-vector multiplications execute directly within memory arrays, eliminating the von Neumann bottleneck that limits conventional digital systems. CIM implementations span diverse technologies: conventional SRAM and DRAM arrays, emerging resistive memories (RRAM, PCM), and organic optoelectronic devices that combine optical sensitivity with synaptic plasticity \citep{xu2025emerging,guo2024organic}.

Organic optoelectronic and electrochemical synaptic devices are particularly attractive for edge vision applications because of their multilevel tunability, low-voltage operation, compatibility with flexible substrates, and inherent photosensitivity that enables direct optical signal processing \citep{xu2025emerging,guo2024organic,photonics2025organicreview}. Recent demonstrations include multilevel memory with analog conductance states \citep{guo2024organic,jung2024organicfilaments}, long retention tails suitable for inference workloads \citep{zeng2023organicmemristor}, and integrated optoelectronic synaptic arrays for visual computing \citep{gebregiorgis2023organiccim,zhang2026opect}.

However, translating these device-level advances into system-level vision performance faces a critical methodological gap. The organic optoelectronic literature remains dominated by device-centric characterizations. Network-level studies are typically limited to small benchmarks or simple perceptrons, while variability, ADC resolution, retention-induced drift, and transfer across fabricated instances are either absent from system analyses or treated only qualitatively \citep{zeng2023organicmemristor,jung2024organicfilaments}. The most relevant precursor for variability-aware organic simulation remains the Monte Carlo study of Alibart \emph{et al.}, which sampled device parameters for a simple organic-memristive perceptron \citep{alibart2016physical}. No clear framework yet connects reported organic-device characteristics to deployment-facing accuracy expectations for modern vision transformers or CNNs.

This gap is particularly acute for hardware-aware training (HAT), which mitigates analog non-idealities by injecting device noise into the forward training path \citep{joshi2020accurate}. Although related in spirit to quantization-aware training \citep{choi2019pact} and to domain randomization in sim-to-real transfer \citep{tobin2017domain}, HAT for CIM faces an additional difficulty: device-to-device (D2D) mismatch is spatially structured and fixed for each hardware instance, unlike independent thermal noise. Standard HAT therefore tends to optimize against one sampled mismatch map. As we show below, this choice can make the model fit a particular hardware instance rather than the deployment distribution, leading to catastrophic accuracy collapse when transferring to fresh hardware.

System-level CIM simulators established effective workflows for inorganic memories. DNN\allowbreak+\allowbreak NeuroSim connected device models to peripheral-circuit costs and network accuracy \citep{peng2020dnnneurosim}. MemTorch embedded memristive non-idealities into PyTorch for end-to-end evaluation \citep{lammie2022memtorch}, AIHWKIT exposed analog HAT and inference simulation in a practical software stack \citep{rasch2021aihwkit}, and CrossSim provides a GPU-accelerated crossbar-accuracy workflow with parasitic resistance, ADC, drift, and lookup-table-based training models \citep{crosssim2026}. These frameworks remain anchored mainly to inorganic resistive memories, and they do not directly capture the coupled photoresponse, retention, and write nonlinearity that define organic optoelectronic inference.

At the same time, the broader organic-device literature provides enough empirical structure to motivate a deployment-facing simulator. Recent array demonstrations report energy-efficient organic device arrays, lithographically patterned near-infrared OPECT arrays, and large-scale optoelectronic synapse arrays for visual neural networks \citep{gebregiorgis2023organiccim,zhang2026opect,zhang2025mooptoelectronic,cui2025multimode}. Related studies also point to active-matrix addressing, spatial optical response control, crosstalk, and thermal stability as practical system constraints rather than purely materials-level observations \citep{visionarch2023crosstalk,amspa2024insensor,jang2023insensor,fuller2020tempresilient,guo2024hightemp}. Taken together, this literature provides plausible ranges for conductance windows, state counts, retention behavior, and array-level artifacts in an organic-CIM workflow.

The mapping problem at the network level is likewise constrained by prior CIM practice. Dense linear operators are natural crossbar targets, whereas depthwise and highly irregular operators often suffer from poor utilization \citep{wu2023bwq,wang2024epim,ge2024allspark}. Recent heterogeneous CIM accelerators for Vision Transformers confirm this pattern: linear projections are mapped to analog arrays, while attention scores, softmax normalization, and dynamic interactions typically remain digital because of their precision requirements and irregular access patterns \citep{kim2025hemlet,lin2024hardsea,bettayeb2024memristorattention}. Low-bit ViT studies further show that attention-adjacent computations remain especially sensitive below 6 bits \citep{liu2021ptqvit,li2022qvit,lin2023vitptq}. These observations motivate a hybrid analog--digital deployment model rather than an all-analog transformer implementation.

Here we present a behavioral simulation framework that bridges this materials-to-system gap. The framework incorporates organic-device-specific non-idealities---photoresponse, retention, and write nonlinearity---through a replaceable device-profile interface and applies a common mixed-signal deployment model to ResNet-18, ConvNeXt-Tiny, and Tiny-ViT-5M. Three constraints dominate in this setting. First, ADC resolution rather than nominal conductance quantization becomes the primary accuracy bottleneck: reducing ADC resolution below 6 bits causes abrupt collapse. Second, standard HAT overfits a fixed D2D realization and fails under fresh-instance transfer; Ensemble HAT, which resamples D2D masks at each training epoch, raises fresh-instance accuracy to $86.37 \pm 1.54\%$. Third, severe nonlinear write ($NL=2.0$) remains difficult to recover under the present gradient-scaling approximation, limiting accuracy to $27.72 \pm 0.82\%$.

The framework is intentionally behavioral rather than circuit-accurate: array-level parasitics (IR drop, sneak paths), thermal effects, and pulse-faithful programming dynamics are not modeled explicitly and are revisited in Section~\ref{sec:discussion}. Its purpose is to identify which device characteristics most constrain edge-vision accuracy before full array hardware becomes available.
